% ---------------------------------------------------------------------------
% Background
%
% The setting and the vocabulary the rest of the paper depends on. Not a
% literature review -- that is 07-related-work.
%
%   - the contest: founded \val{founded}, \val{editions} editions, \val{nsub}
%     submissions, \val{nmatch} matches. Why a public benchmark cannot reproduce
%     it: games original to the venue, and human code spanning a natural skill
%     gradient rather than one reference solution. Do NOT write that original
%     rules are "absent from pretraining data" -- it does not follow (the archive
%     is public) and Section 4 tests it instead. Say the games are the venue's own
%     and point at that test.
%   - the \val{ngames} games, plus the classic games kept as structural controls.
%     Say what they control for, or the word "control" is decoration: they are
%     games whose strategic structure is independently documented and whose
%     K(pi_good) can be argued about without our proxy, so they calibrate the
%     proxy and give the cross-game complexity ordering an external reference
%     point. Without them, every complexity claim rests on a measurement made by
%     the same pipeline it is meant to validate.
%   - the information channel, defined carefully, because the introduction now
%     concedes a limit on it and this is where the concession has to be precise.
%     What a maintained program consumes -- rules, its own source, replays of
%     specific losses -- carries more per sample than a scalar return, but that
%     is a property of the CHANNEL and of the agent reading it, not of the
%     symbolic representation. echo2025hindsight is the demonstration: hindsight
%     rewriting of failed episodes into memory, no program anywhere, same kind of
%     gain. So define the lower bound over "what the loop can read" and the upper
%     over "what the loop can edit", and keep the two definitions separate from
%     the start. Merging them is what produces the unsupportable version of the
%     claim, and it is the version a reader will supply on their own if this
%     section leaves the terms loose.
%   - definitions: heuristic learning; the folk hierarchy; budget, split into
%     learning-only / evaluation / total and kept separate everywhere; the Elo
%     protocol (batch vs online); the three opponent pools (training, frozen
%     test, Elo anchor).
%   - editability, defined against interpretability, since the introduction turns
%     on the distinction and asserts it without room to define it. Interpretability
%     asks whether a human understands the policy; editability asks whether the
%     optimization loop can act on the artifact at all. Programmatic RL already
%     calls its policies human-editable -- cite diprl2026 and only diprl2026,
%     whose abstract says "human-readable and -editable" verbatim. Checked
%     2026-08-18: programmaticrl2024foundations uses none of those words, and
%     llmguidedsearch2024 claims interpretability, which is the very term this
%     definition is distinguishing editability from.
%     Do not claim the word, claim the change in its role: precondition for the
%     optimizer rather than property of the deliverable.
%     The definition should not rest on stipulation, and it does not have to:
%     tig2025think is the case where the two properties visibly come apart. It
%     emits natural-language reasoning for each decision, so it is interpretable
%     on any ordinary reading, while the artifact its learning loop updates is
%     still the weights. Use it as the worked example here; the introduction has
%     room for the citation but not for the anatomy.
%   - K(pi_good), defined here as a concept and measured in Section 3. The one
%     thing this definition must do is separate K(rules) from K(pi*) from
%     K(pi_good); Go is the example that forces the separation, since its rules
%     fit on a page while no compact symbolic policy plays it well. Rule
%     complexity is the wrong x-axis and this is where that is established.
%   - the archive audit protocol: what counts as "no learned component". The
%     criterion has to be decidable -- weight files, training scripts, whether
%     hand-tuned evaluation coefficients count -- because \val{pheur} and
%     \val{prl} are worth exactly as much as this definition is. Classifying the
%     winners is also how RQ3 gets answered, so it is not mere bookkeeping.
% ---------------------------------------------------------------------------

\section{Background}
